\section{Synthetic Torus Validation}
\label{sec:results-torus}

The synthetic torus dataset provides a controlled setting for checking whether
the event diagnostic localizes a known sheet reconfiguration. Unlike the
molecular datasets, the main purpose is validation rather than chemical
interpretation.

\subsection{A Stable Transition Event}

\begin{table}[t]
  \centering
  \caption{Top torus range-space event intervals with threshold \(\theta=0.5\).
  The event scores are symmetric around the controlled transition near
  timestep 50.}
  \label{tab:torus-events}
  \begin{tabular}{lrrrl}
    \toprule
    Interval & Score & Mean best & Weak source/target & Observation \\
    \midrule
    50--51 & 49.08 & 0.246 & 15/15 & strongest transition interval \\
    49--50 & 49.05 & 0.248 & 15/15 & symmetric pre-transition interval \\
    51--52 & 48.09 & 0.296 & 15/15 & symmetric post-transition interval \\
    48--49 & 48.09 & 0.296 & 15/15 & pre-transition shoulder \\
    47--48 & 45.57 & 0.322 & 13/13 & broader transition region \\
    52--53 & 45.57 & 0.322 & 13/13 & broader transition region \\
    \bottomrule
  \end{tabular}
\end{table}

The strongest range-space event occurs at \(50\)--\(51\), with nearly identical
scores on the adjacent intervals \(49\)--\(50\), \(51\)--\(52\), and
\(48\)--\(49\). This symmetric event cluster is the expected behavior for the
synthetic sequence and provides a sanity check for the event diagnostic.

\subsection{Stride And Threshold Stability}

The torus event is stable for all tested thresholds from \(\theta=0.3\) through
\(\theta=0.7\): the top interval remains \(50\)--\(51\). The precomputed stride
views preserve the same transition center. At strides 1, 2, 3, and 4, the top
range intervals are \(50\)--\(51\), \(50\)--\(52\), \(50\)--\(53\), and
\(50\)--\(54\), respectively. Coarser sampling therefore expands the interval
around the same transition instead of moving the detected event elsewhere.

\subsection{Continuing Features}

The tracker also identifies long-lived torus sheet families. Several
threshold-\(0.5\) tracks span the full sequence, including
\(379\rightarrow43\), \(383\rightarrow43\), \(1183\rightarrow43\), and
\(13903\rightarrow43\), each with length 100 and mean continuation score
approximately 0.926. These full-length tracks validate temporal continuity, but
they also illustrate the current non-one-to-one matching model: several tracks
can converge to the same final sheet. We therefore use the torus tracks as
validation evidence, not as a claim of unique one-to-one feature identity.

\subsection{Metric Interpretation}

The combined range score is identical to shape IoU. The torus metric table
shows that bounding-box IoU can have nearly zero mean reference loss on this
synthetic geometry, but its best-target agreement with shape IoU is only 0.598.
We therefore keep shape IoU as the primary correspondence metric and treat the
other shape metrics as secondary diagnostics.

\subsection{Domain Events And Disagreement}

The domain event ranking answers a separate question from the range event
curve. Its strongest intervals are \(7\)--\(8\), \(92\)--\(93\), \(93\)--\(94\),
and \(6\)--\(7\), each with weak domain continuations for all top sheets. These
intervals indicate domain-support behavior in the synthetic sequence and are
not expected to duplicate the range event centered at timestep 50.

Normalized domain/range target disagreement, however, does peak near the
synthetic transition, for example at \(48\)--\(49\) and \(49\)--\(50\). This
provides a useful explanation of the transition: range and domain criteria can
select different targets for the same source sheets near the controlled
reconfiguration. The range-event curve remains the cleanest validation signal,
while the disagreement view explains how the correspondence modes diverge.

\begin{figure*}[t]
  \centering
  \fbox{\parbox{0.92\linewidth}{
  Placeholder: torus validation figure. Show the event-score curve over all
  timesteps, with a highlighted window from 45 to 55 and highlighted intervals
  49--50, 50--51, and 51--52. Include stride annotations showing that strides
  1--4 remain centered at timestep 50.}}
  \caption{The torus event diagnostic localizes the controlled transition near
  timestep 50 and remains stable under coarser sampling.}
  \label{fig:torus-event-validation}
\end{figure*}

\begin{figure}[t]
  \centering
  \fbox{\parbox{0.88\linewidth}{
  Placeholder: torus Sankey transition detail. Show the range-mode Sankey view
  around timesteps 45--55 and highlight one full-length track such as
  \(379\rightarrow43\).}}
  \caption{Torus Sankey detail near the synthetic transition.}
  \label{fig:torus-sankey-detail}
\end{figure}
